\documentclass[a4paper,10pt]{article}

\usepackage[left=2.5cm, right=2cm, top=2cm, bottom=2cm]{geometry}

\usepackage{titlesec}
\titleformat{\section}{\large\bfseries}{\thesection)}{0.5em}{}

\usepackage{enumitem} % Para listas mais compactas
\usepackage{hyperref} % Para links
\usepackage{parskip}  % Para espaçamento entre parágrafos
\usepackage{graphicx}
\usepackage{booktabs}
\usepackage{colortbl}
\usepackage{tabulary}
\usepackage{tabularx}
\usepackage{xcolor}
\usepackage{multirow}

\begin{document}

\begin{center}
\Large
\textbf{SÚMULA CURRICULAR}
\end{center}
\vspace{4mm}

\textsf{\textbf{Rafael Mendonça Duarte}} \\
\textbf{Orcid:} \url{https://orcid.org/0009-0008-7547-0601} \\
\textbf{Currículo Lattes:} \url{http://lattes.cnpq.br/2314499895236009} \\
\textbf{Google Scholar:} \url{https://scholar.google.com.br/citations?hl=pt-BR&user=PZFZg5EAAAAJ}

Rafael Mendonça Duarte, estudante do terceiro período de Sistemas de Informação no Instituto de Ciências Matemáticas e de Computação da Universidade de São Paulo (ICMC-USP), apresenta relevante histórico acadêmico, sem registro de reprovações e com média ponderada de $8,6$. Possui trajetória marcada por premiações em olimpíadas de conhecimento, além de atuação ativa em projetos de pesquisa de Iniciação Científica e extensão universitária. Atualmente, dedica-se ao desenvolvimento e ao aprimoramento de Redes Neurais Convolucionais de Grafos (GCNs) aplicadas à classificação de imagens.


% O docente Lucas Pascotti Valem apresenta produção científica relevante e atuação em projetos que envolvem inteligência artificial, visão computacional e recuperação de imagens. Apresenta incluindo publicações em periódicos e conferências de alto impacto. Possui histórico consistente de excelência acadêmica, evidenciado por diversos prêmios e distinções recebidos ao longo da carreira. Mantém uma ampla rede de colaboração com pesquisadores nacionais e internacionais. %Seus principais interesses de pesquisa incluem recuperação de imagens e aprendizado baseado em grafos.

\section{Formação}
\vspace{-2mm}

\begin{table}[ht!]
    \centering
    % tabularx toma a largura total (\textwidth) e a coluna X expande/quebra linha automaticamente
    \begin{tabularx}{\textwidth}{|X|l|X|}
        \hline
         \textbf{Título} & \textbf{Período} & \textbf{Instituição} \\
         \hline
         Bacharelado em Sistemas de Informação (em andamento) & 2025 – Atual & Instituto de Ciências Matemáticas e Computação (ICMC-USP) \\
         Ensino Médio & 2022 – 2024 & Anglo Alante Jundiaí \\
         \hline
    \end{tabularx}
    %\caption{Caption}
    %\label{tab:my_label}
\end{table}
\vspace{-4mm}

\section{Histórico Profissional e Acadêmico}

\begin{itemize}

\item
\textbf{Cargo:} Aluno de Iniciação Científica. \\
\textbf{Instituição:} ICMC - Instituto de Ciências Matemáticas e de Computação, Universidade de São Paulo (USP). \\
\textbf{Departamento:} Departamento de Ciências de Computação (SCC). \\
\textbf{Remuneração:} Bolsista PRPI, “Programa de Apoio aos Novos Docentes”. \\
\textbf{Período:} Novembro/2025 - Atual. \\
\textbf{Descrição:} Atuo sob orientação do Prof. Dr. Lucas Pascotti Valem no projeto de pesquisa intitulado \textit{"Investigação de Medidas Contextuais de Centralidade para Redes Convolucionais de Grafos em Classificação de Imagens"}. O foco da pesquisa é  investigar a influência de propriedades estruturais do grafo, como medidas de centralidade, nas estratégias de agregação de mensagens em Redes Neurais Convolucionais de Grafos (GCNs) aplicadas à classificação de imagens com poucos dados rotulados.

\item
\textbf{Projeto:} Programa Pró-Aluno. \\
\textbf{Cargo:} Bolsista / Monitor Técnico em Informática. \\
\textbf{Instituição:} ICMC - Instituto de Ciências Matemáticas e de Computação, Universidade de São Paulo (USP). \\
\textbf{Remuneração:} Bolsista Pró-Reitoria de Graduação, “Programa Pró-aluno”. \\ 
\textbf{Período:} Setembro/2025 - Outubro/2025. \\
\textbf{Descrição:} Programa com objetivo de fornecer recursos básicos de informática para os alunos de graduação, para uso exclusivo nas atividades acadêmicas. Atuei no suporte técnico a usuários nos laboratórios de graduação do ICMC, auxiliando no uso de equipamentos, softwares e zelo pela infraestrutura local.

\end{itemize}
\section{Contribuições à Ciência}
\begin{table}[h!]
\centering
\label{tabResPub}
\begin{tabularx}{0.8\textwidth}{Xr} % Ajustado para 2 colunas
    \toprule
    \textbf{Indicador} & \textbf{Quantidade} \\
    \midrule
    \multicolumn{2}{l}{\textit{Produção Científica}} \\
    Artigos Aceitos & 1 \\
    Artigos submetidos & 1 \\
    Artigos em avaliação (\textit{Major Review}) & 1 \\
    \rowcolor{lightgray} \textbf{Total de Artigos} & \textbf{3} \\
    \addlinespace[10pt]
    \multicolumn{2}{l}{\textit{Olimpíadas do Conhecimento}} \\
    Participação & 7 \\
    \rowcolor{lightgray} \textbf{Prêmios e Distinções} & 4 \\
    \bottomrule
\end{tabularx}
\end{table}

\subsection{Artigos Aceitos e Submetidos}
\begin{enumerate}[label=\textbf{\arabic*}.] 

    \item \textbf{DUARTE, R. M.}; PONCIANO, J. R.; VALEM, L. P. \textbf{Gaussian Rank-Based Neighborhood Degree for Graph Neural Networks in Image Classification.} In: \textbf{Proceedings of the Brazilian Symposium on Databases (SBBD 2026)}, São Carlos, SP, Brasil, 2026. \textbf{(Aceito para publicação)}. \\
    $\bullet$ \textbf{Status:} Aceito para publicação em Anais de Congresso (Qualis A4). \\
    $\bullet$ \textbf{DOI (Pré-Impressão):} \url{https://doi.org/10.48550/arXiv.2605.24367}. \\
    $\bullet$ \textbf{Relevância:} A principal contribuição está no método proposto, GRaNDe, uma medida de centralidade dinâmica que integra a análise de vizinhança com a ponderação da distância gaussiana para melhor capturar a importância dos nós no cenário de classificação de imagens, tornando as Redes Neurais de Grafos mais eficazes.

    \item BRITO, G. M.; \textbf{DUARTE, R. M.}; VALEM, L. P. \textbf{Rank Correlation Graphs with Mutual Neighborhood Weighting for GCNs in Image Classification.} In: \textbf{Pattern Recognition Letters}, 2026. \textbf{\textit{(Major Review)}}. \\
    $\bullet$ \textbf{Status:} Submetido / \textit{Major Review} (Qualis A1). \\
    $\bullet$ \textbf{Relevância:} Introduz uma estratégia baseada em vizinhança mútua para ponderação de grafos de correlação. O trabalho traz uma análise comparativa envolvendo diferentes funções de ponderação e topologias, demonstrando a consistência do método proposto para a classificação de imagens.

    % \item \textbf{DUARTE, R. M.}; BRITO, G. M.; VALEM, L. P. \textbf{(\textit{Título Omitido})} In: \textbf{Proceedings of the 39th SIBGRAPI Conference on Graphics, Patterns and Images}, 2026. \textbf{(Submetido)}. \\
    % $\bullet$ \textbf{Status:} Submetido / Sob Avaliação (Qualis A3).

\end{enumerate}

\section{Outras Informações Relevantes}

\subsection{Prêmios e Distinções Acadêmicas}
\begin{enumerate}
    \item \textbf{2024} – Prêmio de Excelência Acadêmica, Anglo Alante Jundiaí.
    \item \textbf{2024} – Medalha de Bronze, Concurso Canguru de Matemática.
    \item \textbf{2024} – Medalha de Bronze, Olimpíada Brasileira de Astronomia e Astronáutica (OBA).
    \item \textbf{2023} – Medalha de Prata, Concurso Canguru de Matemática.
    \item \textbf{2023} – \textit{Position Paper Award}, High School United Nations Human Rights Council 1, FACAMP.
    \item \textbf{2022} – Medalha de Bronze, Mostra Brasileira de Foguetes (MOBFOG).
\end{enumerate}

\subsection{Projetos de Extensão:}
\begin{itemize}

\item
\textbf{Projeto:} Combate à Desinformação Digital: Detecção de Imagens Falsas. \\
\textbf{Cargo:} Integrante. \\
\textbf{Instituição:} ICMC - Instituto de Ciências Matemáticas e de Computação, Universidade de São Paulo (USP). \\
\textbf{Remuneração:} Voluntário. \\
\textbf{Período:} Março/2026 - Atual. \\
\textbf{Descrição:} Atuação como integrante no projeto "Combate à Desinformação Digital: Detecção de Imagens Falsas". Minha função é pesquisar técnicas de aprendizado profundo em visão computacional que corroboram para desenvolvimento e otimização de algoritmos de detecção de imagens falsas.

\textbf{Projeto:} Projeto Egrégora. \\
\textbf{Cargo:} Aluno Coordenador do Projeto e Integrante. \\
\textbf{Instituição:} EESC - Escola de Engenharia de São Carlos, Universidade de São Paulo (USP). \\
\textbf{Remuneração:} Voluntário. \\
\textbf{Período:} Junho/2025 - Atual. \\
\textbf{Descrição:} Projeto focado na área de extensão tecnológica voltado para a recaracterização de dispositivos de hardware do tipo "TV Box". Atuo como coordenador do ramo do projeto voltado para a montagem de clusters de computação distribuída de baixo custo e implantação de servidores locais de banco de dados.

\item
\textbf{Projeto:} Curso de Informática para a Terceira Idade - 3ª Edição. \\
\textbf{Cargo:} Aluno Monitor. \\
\textbf{Instituição:} ICMC - Instituto de Ciências Matemáticas e de Computação, Universidade de São Paulo (USP). \\
\textbf{Remuneração:} Voluntário. \\
\textbf{Período:} Março/2026 - Junho/2026. \\
\textbf{Descrição:} Atividade de extensão voltada à inclusão digital da comunidade externa de São Carlos. Atuação como aluno monitor, auxiliando no esclarecimento de dúvidas e acompanhamento individual dos alunos.
\end{itemize}




% ----------------------------------------------------

% \section{Financiamentos à Pesquisa}

% Principais financiamentos recentes (últimos 5 anos):

% \begin{itemize}
%     \item \textbf{FAPESP APR - Aprendizado Contextual de Similaridade Aplicado a Redes Convolucionais de Grafos para Classificação e Recuperação de Imagens}
%     \begin{itemize}
%     \item \textbf{Período:} 02/2026 -- 01/2029 (atual).
%     \item \textbf{Processo:} 	2025/10602-5.
%         \item \textbf{Descrição:} O projeto investiga o uso de Redes Convolucionais de Grafos (GCNs) para tarefas de classificação e recuperação de imagens, com foco em cenários com poucos dados rotulados. A proposta explora informações contextuais e de vizinhança para aprimorar diferentes etapas do processo, como a construção de grafos, estratégias de treinamento e aprendizado de representações.
%     \end{itemize}

%     \item \textbf{Programa de Apoio aos Novos Docentes USP – 2025}
%     \begin{itemize}
%     \item \textbf{Período:} 04/2025 -- 04/2027 (atual).
%     \item \textbf{Portaria:} PRPI nº 1032, de 17 de março de 2025 (DOE-SP 18/03/2025).
%         \item \textbf{Descrição:} Oferece recursos para apoiar atividades iniciais de docentes recém-admitidos.
%     \end{itemize}

%     % \item
%     % \textbf{Bolsa de Doutorado Sanduíche Fulbright – Estágio no Exterior}
%     % \begin{itemize}
%     %     \item \textbf{Título:} \textit{Graph Convolutional Networks for Weakly Supervised Representation Learning on Multimedia Classification and Retrieval}.
%     %     \item \textbf{Universidade:} \textit{Temple University, Philadelphia - PA, USA}.
%     %     \item \textbf{Supervisor:} Prof. Dr. Longin Jan Latecki.
%     %     \item \textbf{Período:} 01 de setembro de 2022 - 31 de maio de 2023 (9 meses).
%     %     \item \textbf{Resumo das atividades}: O projeto de pesquisa teve como objetivo explorar o uso de redes neurais convolucionais baseadas em grafos (GCNs) para aplicações em cenários de recuperação e classificação de imagens com poucos ou nenhum dado rotulado para treinamento.
%     % \end{itemize}

%     % \item
%     % \textbf{Bolsista de Treinamento Técnico FAPESP (Nível IV)}
%     % \begin{itemize}
%     %     \item \textbf{Título:} Suporte para ambiente computacional e execução de experimentos: aprendizado fracamente supervisionado e fusão de métodos de classificação.
%     %     \item \textbf{Universidade:} Universidade Estadual Paulista Júlio de Mesquita Filho, UNESP, Rio Claro.
%     %     \item \textbf{Supervisor:} Prof. Dr. Daniel Carlos Guimarães Pedronette.
%     %     \item \textbf{Período:} 01 de novembro de 2020 - 30 de junho de 2021 (8 meses).
%     %     \item \textbf{Processo:} 2020/11366-0.
%     %     \item \textbf{Convênio/Acordo:} Microsoft Research.
%     %     \item \textbf{Resumo das atividades}: Técnicas de aprendizado de máquina geralmente dependem de grandes conjuntos de dados rotulados para construir modelos preditivos. No entanto, a limitação de conjuntos de dados rotulados motiva a investigação de abordagens fracamente supervisionadas. Foram desenvolvidos métodos fracamente supervisionados para busca e classificação de imagens.
%     % \end{itemize}

% \end{itemize}

% Outros projetos em que participo:

% \begin{itemize}
%     \item \textbf{Projeto FAPESP PCD/CCD – IntegraSanca: Conectando Dados, Transformando São Carlos (2025/07160-0):} visa investigar a integração e interoperabilidade de dados da prefeitura de São Carlos, atuando nas áreas de banco de dados, inteligência artificial, segurança da informação e engenharia de software, com o objetivo de modernizar a gestão e alinhar o município ao conceito de cidades inteligentes.

%     \item \textbf{Projeto FAPESP Temático – Gestão Inteligente de Dados Multimodais de Saúde para Tomada de Decisão em Cenários de Big Data - IHealth-MD  (2024/13328-9):} visa desenvolver algoritmos de Big Data e IA para integrar e extrair conhecimento de dados de saúde multimodais oriundos de registros eletrônicos de pacientes, com foco em doenças pulmonares e cardiovasculares. 

%     \item \textbf{Projeto FAPESP – NLP com Transformers (2024/04890-5):} Parceria com UNESP e Idiap/Suíça, no âmbito de cooperação FAPESP–SNSF. Foco em recuperação aumentada e verificação de fatos com LLMs. Participação como pesquisador (não bolsista).
% \end{itemize}

% \subsection{Rede de Colaboradores}

% O docente mantém uma ampla rede de colaboração com pesquisadores no Brasil e no exterior. A tabela a seguir apresenta os principais professores que atuaram como coautores de publicações conjuntas.

% %O docente possui uma ampla rede de colaboradores nacionais e internacionais. A tabela apresenta todos os professores pesquisadores que são coautores do docente.

% \begin{table}[ht!]
% \centering
% \begin{tabular}{c|p{4.6cm}|p{6.5cm}}
% \cline{2-3}
%  & \textbf{Colaborador} & \textbf{Instituição}\\
% \hline
%  & Daniel C. G. Pedronette & \multirow{4}{*}{UNESP, Rio Claro, SP.} \\
% \cline{2-2}
%  & Ivan Rizzo Guilherme &  \\
% \cline{2-2}
%  & Alexandro Baldassin &  \\
% \cline{2-2}
%  & Fabricio Breve &  \\
% \cline{2-2}
% % & Denis H. P. Salvadeo &  \\
% \cline{2-3}
%  & João Paulo Papa & UNESP, Bauru, SP. \\
% \cline{2-3}
%  & Renato Tinós & \multirow{3}{*}{USP, Ribeirão Preto, SP.} \\
% \cline{2-2}
% \textbf{Nacionais} & Liang Zhao &  \\
% \cline{2-2}
%  & Luiz Otávio Murta Junior &  \\
% \cline{2-3}
%  & Roseli A. Francelin Romero & \multirow{2}{*}{USP, São Carlos, SP.} \\
% \cline{2-2}
%  & Mirela T. Cazzolato &  \\
% \cline{2-3}
%  & Jurandy Almeida & UFSCAR, Sorocaba, SP. \\
% \cline{2-3}
%  & Edson Borin & UNICAMP, Campinas, SP.\\
% \cline{2-3}
%  & Everton C. Tetila & UFGD, Dourados, MS. \\
% \hline
%  & Longin Jan Latecki &  Temple University, Filadélfia,\newline PA, E.U.A. \\
% \cline{2-3}
%  & Mohand Said Allili & Université du Quebec en Outaouais, Gatineau, Canadá.  \\
% \cline{2-3}
% \textbf{Internacionais} & Ricardo da Silva Torres & Wageningen University and \newline Research, Wageningen, Gueldres, Holanda. \\
% \cline{2-3}
%  & Fabio Augusto Faria & IST/ULisboa, Lisboa, Portugal. \\
% \cline{2-3}
%  & Theodore Gyle Lewis & Naval Postgraduate School, Monterey, CA, E.U.A. \\
% \cline{2-3}
%  & Andre Freitas & University of Manchester, England, UK. \\
% \hline
% \end{tabular}
% \label{tab:principais_profpes}
% \end{table}

% % \begin{table}[ht!]
% % \centering
% % %\caption{Professores pesquisadores que colaboram ou colaboraram com o candidato.}
% % \begin{tabular}{c|p{4.6cm}|p{6.5cm}}
% % \cline{2-3}
% %  & \textbf{Colaborador} & \textbf{Instituição}\\
% % \hline
% %  & Daniel C. G. Pedronette &  UNESP, Rio Claro, SP. \\
% % \cline{2-3}
% % & Ivan Rizzo Guilherme & UNESP, Rio Claro, SP. \\
% % \cline{2-3}
% % & Alexandro Baldassin & UNESP, Rio Claro, SP. \\
% % \cline{2-3}
% % & Fabricio Breve & UNESP, Rio Claro, SP. \\
% % \cline{2-3}
% % & João Paulo Papa & UNESP, Bauru, SP.  \\
% % \cline{2-3}
% % & Renato Tinós & USP, Ribeirão Preto, SP. \\
% % \cline{2-3}
% % \textbf{Nacionais} & Liang Zhao & USP, Ribeirão Preto, SP. \\
% % \cline{2-3}
% % & Luiz Otávio Murta Junior & USP, Ribeirão Preto, SP. \\
% % \cline{2-3}
% % & Roseli A. Francelin Romero & USP, São Carlos, SP. \\
% % \cline{2-3}
% % & Mirela T. Cazzolato & USP, São Carlos, SP. \\
% % \cline{2-3}
% % & Jurandy Almeida & UFSCAR, Sorocaba, SP. \\
% % \cline{2-3}
% % & Edson Borin & UNICAMP, Campinas, SP.\\
% % \cline{2-3}
% % & Everton C. Tetila & UFGD, Dourados, MS. \\
% % \hline
% %  & Longin Jan Latecki &  Temple University, Filadélfia,\newline PA, E.U.A. \\
% % \cline{2-3}
% % & Mohand Said Allili & Université du Quebec en Outaouais, Gatineau, Canadá.  \\ 
% % \cline{2-3}
% % \textbf{Internacionais} & Ricardo da Silva Torres & Wageningen University and \newline  Research, Wageningen, Gueldres, Holanda. \\
% %  \cline{2-3}
% %   & Fabio Augusto Faria & IST/ULisboa, Lisboa, Portugal. \\
% % \cline{2-3}
% % & Theodore Gyle Lewis & Naval Postgraduate School,
% % Monterey, CA, E.U.A. \\
% % \cline{2-3}
% % & Andre Freitas & University of Manchester, England, UK \\
% % \hline
% % \end{tabular}
% % \label{tab:principais_profpes}
% % \end{table}

% % \begin{table}[ht!]
% % \centering
% % %\caption{Professores pesquisadores que colaboram ou colaboraram com o candidato.}
% % \begin{tabular}{c|p{4.6cm}|p{6.5cm}}
% % \cline{2-3}
% %  & \textbf{Colaborador} & \textbf{Instituição}\\
% % \hline
% %  & Daniel C. G. Pedronette &  UNESP, Rio Claro, SP, Brasil.  \\
% % \cline{2-3}
% % & Jurandy Almeida & UFSCAR, Sorocaba, SP, Brasil. \\
% % \cline{2-3}
% %  & João Paulo Papa & UNESP, Bauru, SP, Brasil.  \\
% % \cline{2-3}
% % \cline{2-3}
% % \textbf{Nacionais} & Ivan Rizzo Guilherme & UNESP, Rio Claro, SP, Brasil. \\
% % \cline{2-3}
% % & Edson Borin & UNICAMP, Campinas, SP, Brasil.\\
% % \cline{2-3}
% % & Alexandro Baldassin & UNESP, Rio Claro, SP, Brasil. \\
% % \cline{2-3}
% % & Fabricio Breve & UNESP, Rio Claro, SP, Brasil. \\
% % \hline
% %  & Longin Jan Latecki &  Temple University, Filadélfia,\newline PA, E.U.A. \\
% % \cline{2-3}
% % & Mohand Said Allili & Université du Quebec en Outaouais, Gatineau, Canadá.  \\ 
% % \cline{2-3}
% % \textbf{Internacionais} & Ricardo da Silva Torres & Wageningen University and \newline  Research, Wageningen, Gueldres, Holanda. \\
% %  \cline{2-3}
% %   & Fabio Augusto Faria & IST/ULisboa, Lisboa, Portugal. \\
% % \cline{2-3}
% % & Theodore Gyle Lewis & Naval Postgraduate School,
% % Monterey, CA, E.U.A. \\
% % \cline{2-3}
% % & Andre Freitas & University of Manchester, England, UK \\
% % \hline
% % \end{tabular}
% % \label{tab:principais_profpes}
% % \end{table}


% %\vspace{-2mm}
% \subsection{Principais Softwares Publicados}

% %Com a finalidade de divulgar e facilitar a reprodutibilidade dos trabalhos publicados, diversos \textit{frameworks} de \textit{software} foram disponibilizados.

% \begin{itemize}
%     \item
%     \textbf{Unsupervised Distance Learning Framework (UDLF)}
%     \begin{itemize}
%         \item Disponibiliza 11 métodos de aprendizado não supervisionado para pós-processamento de resultados de busca de imagens. Disponível em \url{https://github.com/UDLF/UDLF}.
%     \end{itemize}

%     \item \textbf{Demais softwares disponíveis em:} \url{https://lucasvalem.com/softwares}.
%     %\textbf{Python Wrapper for Unsupervised Distance Learning Framework (pyUDLF)}
%     %\begin{itemize}
%     %    \item \textbf{Descrição:} O pyUDLF é um \textit{wrapper} em Python para o UDLF, que proporciona a execução do mesmo de maneira mais simplificada em uma linguagem de mais alto nível. A linguagem Python é amplamente utilizada por \textit{frameworks} de aprendizado de máquina e processamento de imagens. \\
%     %    \textbf{\mbox{*}} Disponível em \url{https://github.com/UDLF/pyUDLF}.
%     %\end{itemize}

% \end{itemize}

\end{document}